% 4-page anonymized workshop submission.
% Compile with pdflatex. Requires fig_main.pdf in the same directory
% (or adjust \graphicspath).
% This is a COMPRESSED, ANONYMIZED derivative of the full 8-page paper.
% Do not edit the full version from here.
\documentclass[10pt]{article}
\usepackage[margin=1in]{geometry}
\usepackage{amsmath,amssymb}
\usepackage{booktabs}
\usepackage{graphicx}
\graphicspath{{../figures/}{./}}
\usepackage[numbers,sort&compress]{natbib}
\usepackage[hidelinks]{hyperref}
\usepackage{titlesec}
\usepackage{multirow}

% tighten spacing to fit 4 pages
\titlespacing*{\section}{0pt}{1.4ex plus 0.3ex}{0.8ex}
\titlespacing*{\subsection}{0pt}{1.1ex plus 0.2ex}{0.5ex}
\setlength{\textfloatsep}{8pt plus 2pt minus 2pt}
\setlength{\intextsep}{8pt plus 2pt minus 2pt}
\setlength{\abovecaptionskip}{4pt}

\newcommand{\rinse}{\textsc{rinse}}

\title{\vspace{-3em}Episode Length Is a Hidden Design Choice:\\
Demonstration-Curation Metrics Invert Once You Control For It}

\author{Anonymous Author(s)\\ \small Affiliation withheld for review}
\date{}

\begin{document}
\maketitle
\vspace{-2em}

\begin{abstract}
\noindent
Trajectory-smoothness scores are increasingly used to curate
imitation-learning demonstrations, on the premise that smoother
demonstrations are better ones. We show that on real robomimic data these
scores are largely measuring \emph{how long the episode took}, and that the
direction of their empirical claim flips once episode duration is held
fixed. On can/mh, Spectral Arc Length separates better from worse
teleoperators with AUROC 0.952, but episode length alone achieves 0.957,
and inside a duration-matched band SAL falls to 0.628 while raw mean jerk,
the crude baseline such metrics are motivated against, rises to 0.896. The
pattern replicates without retuning on square/mh (SAL 0.898$\to$0.480, jerk
0.431$\to$0.896, length control 0.936$\to$0.407). On a fixed-horizon split
where duration cannot carry information, both published metrics sit at or
below chance against task success, and curating training data by the
contact-aware one costs a third of downstream policy success relative to
random selection of the same size (10 seeds, $p{=}0.0001$ after correction
for multiple comparisons). Under a diffusion policy that ordering holds and
jerk inverts again, beating both baselines. Whether a duration control is
applied is not a robustness detail: it decides the sign of the reported
result. We argue it belongs in the minimum reporting set for any
demonstration-scoring claim.
\end{abstract}

\section{The design choice underneath the claim}

A demonstration-curation metric is evaluated by showing that it ranks known
good demonstrations above known bad ones. That evaluation contains a design
choice that is almost never reported: whether the good and bad
demonstrations differ in \emph{duration}. When they do, and in human
teleoperation data they systematically do, because worse operators take
longer, any statistic that grows with trajectory length inherits an
apparent discriminative power that has nothing to do with what the metric
was designed to measure.

This is not a hypothetical failure mode. Prior audit work on a
synthetic-defect testbed found that most curation metrics examined achieved
near-perfect detection scores by proxying episode
length~\citep{bedi-policy}, and controlled for it by truncating all
trajectories to a fixed prefix. We ask what happens to \emph{named,
published} metrics on \emph{real} operator labels when the same control is
applied, and find that the control does not merely shrink the effect. It
reverses the ranking. The metric that looks best without the control is
near chance with it; the crude baseline it was introduced to supersede
becomes the strongest detector.

We report three things. (1) On two robomimic tasks, smoothness-metric
performance tracks the episode-length baseline almost exactly, and collapses
under duration matching. (2) On a fixed-horizon split where duration is
constant by construction, both metrics of a recent published method are at
or below chance against task success, with the contact-aware metric
inverted. (3) Filtering training data by these scores produces policies
worse than no filtering at all. The substantive contribution is not that
these particular metrics are weak; it is that a single unreported
preprocessing decision determines the sign of the conclusion, which makes it
a reporting requirement rather than a robustness check.

\textbf{Related work.} \rinse{}~\citep{rinse} and PSD~\citep{psd} are recent
smoothness- and spectral-power-based curation methods; both validate on
success-filtered pools, and neither evaluates on data containing failures.
The audit line we build on~\citep{bedi-audit,bedi-policy,bedi-phase}
established the detection-vs-curation decoupling and documented the
length confound, but on generic metric implementations, synthetically
injected defects, and with length controlled by truncation. We extend it to
named published metrics, real operator and success labels, and a
fixed-horizon split where the confound is excluded by construction rather
than by discarding signal. Separately,~\citep{empirical-science,roboeval}
argue that success rate under-describes policy quality
and~\citep{benchmarking-audit} shows benchmarks admit shortcut solutions;
our result is the data-side complement, where the scores used to
\emph{select} training data are themselves shortcut-driven.

\section{Setup}

\textbf{Metrics.} Spectral Arc Length (SAL) and Trajectory-Envelope Distance
(TED) as specified by \rinse{}~\citep{rinse}, a recent smoothness-based
curation method; TED is reimplemented from its appendix, including the
gripper-contact partition that is its distinguishing feature. Baselines:
mean jerk (third-difference magnitude of end-effector position) and
end-effector path length. All are computed from proprioception at 20\,Hz and
returned as badness scores. We also report the per-timestep normalization
SAL$/T$, since SAL's arc length sums over $T/2$ spectral bins and therefore
grows with duration by construction.

\textbf{Data.} Public robomimic~\citep{robomimic} low-dim. \emph{can/mh} and
\emph{square/mh}: 300 successful human demonstrations each with released
operator skill tiers (better/okay/worse, 100 each); durations 98--1050 and
123--1051 steps. \emph{can/mg}: 3900 scripted demonstrations, 18.4\%
successful, every episode exactly 150 steps. Can is the only robomimic task
shipping failure-containing splits, so the success-label analysis is
confined to it by benchmark availability rather than by choice.

\textbf{Controls.} Two, applied independently. \emph{Normalization}: divide
by $T{-}1$. \emph{Duration matching}: restrict to episode-length bands of
increasing tightness, and report the AUROC of episode length itself in each
band: a band is only a control where that number is near 0.5. All AUROCs
carry 95\% bootstrap CIs (2000 resamples); full tables are in the
supplementary material.

\section{The confound, and what controlling for it does}

\begin{figure}[t]
\centering
\includegraphics[width=0.92\linewidth]{fig_main.pdf}
\caption{\textbf{Smoothness tracks duration, not quality.} (a) Skill AUROC
on can/mh across tightening duration bands. SAL (red) follows the
episode-length baseline (gray, dashed) down; jerk (blue) rises as the
confound is removed; the two cross. In the tightest band length itself is at
chance (0.485), confirming the control. (b) On the fixed-horizon can/mg
split, duration carries no information by construction; against task
success, jerk retains signal while SAL is at chance and TED is inverted.}
\label{fig:main}
\end{figure}

\textbf{Uncontrolled, SAL looks excellent, and so does a stopwatch.} On
can/mh, SAL separates better from worse operators at AUROC 0.952 [0.922,
0.975]. Episode length alone achieves 0.957 [0.929, 0.980]. SAL's
correlation with $T$ is $+0.97$, and its tier means (100.4 / 57.5 / 44.9)
track the tier durations (304 / 181 / 143 steps) almost linearly.

\textbf{Both controls remove it.} Normalization drops SAL to 0.582 [0.497,
0.662] and its length correlation to $+0.07$. Duration matching does the
same thing more directly: across four bands SAL falls monotonically while
jerk rises monotonically, and they cross (Table~\ref{tab:both},
Figure~\ref{fig:main}a). In the band where episode length is itself at
chance, SAL's CI includes 0.5 and jerk reaches 0.896.

\textbf{It replicates on a second task, unchanged.} Applying the identical
pipeline to square/mh, with no retuning: SAL's duration correlation is
$+0.96$, its uncontrolled AUROC 0.898 against a length baseline of 0.936,
and in the controlled band SAL sits at 0.480 while jerk reaches 0.896,
the same value it reaches in Can's tightest band. Two task-specific
differences are worth stating rather than smoothing over: jerk is
\emph{inverted} on Square before any control (0.431, versus 0.584 on Can),
and path length stays strong in Square's controlled band where it collapsed
on Can. The uncontrolled ranking of metrics is task-dependent; the fact that
SAL's advantage does not survive the control is not.

\begin{table}[t]
\centering
\small
\begin{tabular}{llcccc}
\toprule
& & \multicolumn{4}{c}{duration band, loosest $\to$ tightest} \\
\cmidrule(l){3-6}
& & full & band 2 & band 3 & controlled \\
\midrule
\multirow{4}{*}{\rotatebox{90}{can/mh}}
& SAL        & 0.952 & 0.793 & 0.759 & 0.628 \\
& jerk       & 0.584 & 0.862 & 0.906 & \textbf{0.896} \\
& TED        & 0.558 & 0.653 & 0.631 & 0.359 \\
& length $T$ & 0.957 & 0.801 & 0.753 & \textbf{0.485} \\
\midrule
\multirow{4}{*}{\rotatebox{90}{square/mh}}
& SAL        & 0.898 & 0.759 & 0.529 & 0.480 \\
& jerk       & 0.431 & 0.570 & 0.782 & \textbf{0.896} \\
& TED        & 0.663 & 0.685 & 0.599 & 0.557 \\
& length $T$ & 0.936 & 0.824 & 0.613 & \textbf{0.407} \\
\bottomrule
\end{tabular}
\caption{Better-vs-worse operator AUROC as the duration band tightens, on
two tasks. The \emph{length} row is the control: only the rightmost column
is duration-matched. SAL falls and jerk rises in lockstep with it, on both
tasks. Bootstrap CIs, exact band edges, and sample sizes are in the
supplementary material; the tightest bands are small ($n{=}40$ and $n{=}30$)
and the claim rests on the monotone trend rather than any single cell.}
\label{tab:both}
\end{table}

\textbf{With duration removed by construction, the metrics fail outright.}
The can/mg split fixes every episode at 150 steps, so no metric can exploit
duration and the length baseline is exactly 0.500. Against the real
task-success label, jerk reaches 0.706, SAL 0.454, path length 0.421, and
TED 0.376 (Figure~\ref{fig:main}b). TED's inversion means filtering by TED
\emph{preferentially retains failed demonstrations}.

\section{Does it matter downstream? Yes, in the harmful direction}

Detection AUROC is not what practitioners care about, so we filter can/mg by
each metric at a 50\% keep fraction, hold the training pool size fixed at
1950 demonstrations, and train GMM behavior-cloning policies (10 seeds each,
50 rollouts per evaluation, best checkpoint per seed). Mean best success
rates: no curation 0.380, random 50\% subset 0.360, SAL 0.332, jerk 0.332,
path length 0.310, TED 0.240. An oracle keeping only true successes reaches
0.892, roughly $2.5\times$ the best metric, so the signal is
unambiguously present in the pool and the metrics are not finding it.

Testing those gaps properly (difference-of-means bootstrap, 20{,}000
resamples, plus Welch two-sample tests, Bonferroni-corrected across eight
comparisons at $\alpha = 0.006$): \textbf{TED} is below random by $-0.120$
[$-0.166$, $-0.076$], $p = 0.0001$, and below no-curation by $-0.140$,
$p < 0.0001$; both survive correction. Path length is directionally harmful
with weaker support ($p = 0.037$, $p = 0.0071$), clearing conventional
significance but not correction. SAL and jerk are below both baselines in
the mean but not distinguishable from them ($p \geq 0.06$). We report them
that way rather than as evidence of harm: an earlier 5-seed pass appeared to
show all four metrics underperforming, and that pattern did not survive the
seed expansion. Since this paper's argument is that under-powered evaluation
yields confident claims in the wrong direction, we hold our own results to
the same standard. Random and no-curation are themselves indistinguishable
($\Delta = -0.020$, CI spanning zero), so a metric's deficit is attributable
to \emph{which} demonstrations it discards, not to pool size. Full grid and
per-comparison tests in supplementary.

\paragraph{Across architectures.} A diffusion-policy replication (five
conditions, five seeds, matched epoch and evaluation budget; supplementary)
inverts one metric and leaves the other in place. TED is worst under both
(0.240 BC, 0.092 diffusion), which is what one expects if the deficit
belongs to the metric; but at five seeds the grid resolves only differences
above 0.054, and TED's deficit there ($-0.028$, CI $[-0.052, +0.000]$) falls
below it, so the \emph{ordering} replicates while the \emph{harm} stays
established under BC only. Jerk reverses sign, beating both baselines by
$+0.096$ and $+0.092$ ($p = 0.0007$, $0.0011$, surviving correction) after
being inert under BC: the crude baseline has curation value under the right
policy class.

Qualitative inspection explains the direction (Figure~\ref{fig:qual}). On
this task successful manipulation requires sustained high-deviation motion
through the contact phase, while failure looks like smooth, low-amplitude
hovering, a robot that lifts and never commits to a placement. A metric
rewarding smoothness is measuring commitment with the sign flipped.

\begin{figure}[!ht]
\centering
\includegraphics[width=0.88\linewidth]{fig_qual_pair.pdf}
\caption{\textbf{Failure is smooth; success is not.} A successful (left) and
failed (right) can/mg demonstration, with TED's contact-segment boundaries
(red), per-segment B\'ezier envelope (dashed), and envelope residual
(bottom). The success carries one high-deviation segment spanning timesteps
12--96, coinciding with the descend-and-grasp. In the failure the only
comparable deviation is the free-space approach (0--24); the remaining 126
steps are a flat tail of repeated gripper cycles that never complete a
placement. TED averages over the alignment path and so rates the failure as
the cleaner trajectory (0.042 vs.\ 0.144). Illustrative of the mechanism;
the quantitative claim is Figure~\ref{fig:main}b.}
\label{fig:qual}
\end{figure}

\section{What should be reported}

The episode-length confound is a substrate effect in a precise sense: it is
a property of how the evaluation data was assembled, invisible in the
metric's definition, that determines whether the reported effect points up
or down. It is cheap to check and cheap to report, and we suggest three
things belong in any demonstration-scoring result.

\textbf{Report the length baseline.} Episode length alone, scored as if it
were a curation metric, against the same label. If it matches the proposed
metric, the proposed metric has not yet been shown to do anything else. This
one number would have flagged the issue in both tasks we examined.

\textbf{Report a duration-controlled condition, and report its control.} A
duration-matched band or a normalization, \emph{with} the length AUROC
inside that condition, so a reader can verify the control worked. A band
where length still scores 0.82 is not a control, and we report our own
intermediate bands partly to make that visible.

\textbf{Separate the label from the deployment target.} Skill tiers and task
success are different labels and these metrics behave oppositely on them.
Curation methods are validated on success-filtered pools and deployed to
remove failures from mixed pools; that gap is where the inversion lives.

\section{Limitations}

Detection results cover two robomimic tasks; the success-label and training
results cover one, because Can is the only robomimic task with
failure-containing splits, and extending them requires a different benchmark
such as LIBERO. Two policy classes downstream (GMM behavior cloning,
diffusion); detection results are policy-free. The diffusion grid covers two
of the four metrics at five seeds, so TED's cross-architecture claim is
about ordering, not harm; its absolute rates stay $3\times$ below BC at a
matched epoch and evaluation budget, which rules out the under-convergence
explanation we gave for the same gap in an earlier pilot and leaves it
unexplained. Corrections are applied within each architecture; pooling all
twelve tests ($\alpha = 0.00417$) changes no conclusion. TED is our
reimplementation from the published
appendix, not the authors' code; it agrees broadly with SAL's ranking on
clean data as a sanity check, but a discrepancy cannot be excluded. The
tightest duration bands are small ($n{=}40$, $n{=}30$), and we treat them as
the endpoint of a trend rather than as standalone evidence. Failures in
can/mg are machine-generated; human failure modes may differ, and a
human-failure pool is the clearest next step.

\paragraph{Use of AI assistance.} An AI assistant was used for code
scaffolding, LaTeX drafting, and editing support. All experiments were
designed, executed, and verified by the author(s); all reported numbers were
produced by the released code and checked against raw outputs.

{\small
\bibliographystyle{plainnat}
\begin{thebibliography}{9}
\bibitem[Kulkarni et al.(2026)]{rinse}
S.~Kulkarni, R.~Dhar, Y.~Cui.
\newblock Learning from the Best: Smoothness-Driven Metrics for Data Quality
in Imitation Learning. arXiv:2604.23000, 2026.

\bibitem[Bedi(2026a)]{bedi-audit}
A.~Bedi.
\newblock Auditing Demonstration Curation Metrics: Action-Only Scorers Fail
on the Structural Defects That Degrade Imitation Policies.
arXiv:2606.05588, 2026.

\bibitem[Bedi(2026b)]{bedi-policy}
A.~Bedi.
\newblock What Demonstration Curation Metrics Do to Your Policy.
arXiv:2606.10229, 2026.

\bibitem[Bedi(2026c)]{bedi-phase}
A.~Bedi.
\newblock Phase-Localized Curation Does Not Help: A Negative Result on
Per-Phase Metric Selection for Demonstration Filtering. arXiv:2606.15064,
2026.

\bibitem[Sojib and Begum(2026)]{psd}
N.~Sojib, M.~Begum.
\newblock An Efficient Metric for Data Quality Measurement in Imitation
Learning. arXiv:2605.01544, 2026.

\bibitem[Mandlekar et al.(2021)]{robomimic}
A.~Mandlekar, D.~Xu, J.~Wong, S.~Nasiriany, C.~Wang, R.~Kulkarni,
L.~Fei-Fei, S.~Savarese, Y.~Zhu, R.~Mart\'in-Mart\'in.
\newblock What Matters in Learning from Offline Human Demonstrations for
Robot Manipulation. CoRL, 2021.

\bibitem[Kress-Gazit et al.(2024)]{empirical-science}
H.~Kress-Gazit, K.~Hashimoto, N.~Kuppuswamy, P.~Shah, P.~Horgan, G.~Richardson,
S.~Feng, B.~Burchfiel.
\newblock Robot Learning as an Empirical Science: Best Practices for Policy
Evaluation. arXiv:2409.09491, 2024.

\bibitem[Wang et al.(2026)]{roboeval}
Y.~R.~Wang, C.~Ung, C.~Tan, G.~Tannert, J.~Duan, J.~Li, A.~Le, R.~Oswal,
M.~Grotz, W.~Pumacay, Y.~Deng, R.~Krishna, D.~Fox, S.~Srinivasa.
\newblock RoboEval: Where Robotic Manipulation Meets Structured and Scalable
Evaluation. arXiv:2507.00435, 2026.

\bibitem[Jiang et al.(2026)]{benchmarking-audit}
T.~Jiang, X.~Tan, S.~Wheeler, L.~Sun, T.~W.~Ayalew, M.~Walter.
\newblock What Are We Actually Benchmarking in Robot Manipulation?
arXiv:2606.04233, 2026.
\end{thebibliography}
}

\clearpage
\appendix
\begin{center}\Large\textbf{Supplementary Material}\end{center}
\vspace{1em}

\section{Full duration-band tables with bootstrap CIs}

\begin{table}[h]
\centering
\small
\begin{tabular}{lcccc}
\toprule
& full ($n{=}200$) & $T \in [120,220]$ ($n{=}98$) & $T \in [130,190]$ ($n{=}74$) & $T \in [140,175]$ ($n{=}40$) \\
\midrule
SAL       & 0.952 [.922,.975] & 0.793 [.690,.881] & 0.759 [.628,.875] & 0.628 [.378,.856] \\
SAL$/T$   & 0.582 [.497,.662] & 0.600 [.479,.715] & 0.603 [.446,.755] & 0.658 [.342,.923] \\
TED       & 0.558 [.480,.636] & 0.653 [.522,.776] & 0.631 [.468,.779] & 0.359 [.147,.583] \\
jerk      & 0.584 [.505,.665] & 0.862 [.783,.928] & 0.906 [.832,.965] & 0.896 [.783,.980] \\
path len. & 0.783 [.718,.842] & 0.590 [.449,.730] & 0.624 [.461,.781] & 0.429 [.215,.645] \\
\midrule
length $T$ & 0.957 [.929,.980] & 0.801 [.689,.901] & 0.753 [.594,.891] & \textbf{0.485} [.227,.740] \\
\bottomrule
\end{tabular}
\caption{can/mh, better-vs-worse operator AUROC, 95\% bootstrap CIs (2000
resamples). The tightest band retains 33 better and 7 worse
demonstrations.}
\end{table}

\begin{table}[h]
\centering
\small
\begin{tabular}{lcccc}
\toprule
& full ($n{=}200$) & $T \in [150,300]$ ($n{=}111$) & $T \in [180,260]$ ($n{=}55$) & $T \in [200,240]$ ($n{=}30$) \\
\midrule
SAL       & 0.898 [.852,.936] & 0.759 [.658,.849] & 0.529 [.344,.702] & 0.480 [.250,.701] \\
SAL$/T$   & 0.458 [.376,.536] & 0.485 [.372,.602] & 0.468 [.300,.653] & 0.520 [.280,.738] \\
TED       & 0.663 [.584,.736] & 0.685 [.577,.788] & 0.599 [.427,.758] & 0.557 [.335,.785] \\
jerk      & 0.431 [.349,.509] & 0.570 [.455,.687] & 0.782 [.618,.923] & 0.896 [.724,1.000] \\
path len. & 0.931 [.893,.963] & 0.817 [.726,.895] & 0.796 [.654,.913] & 0.887 [.742,.991] \\
\midrule
length $T$ & 0.936 [.902,.964] & 0.824 [.738,.898] & 0.613 [.456,.766] & \textbf{0.407} [.197,.629] \\
\bottomrule
\end{tabular}
\caption{square/mh, same format. The tightest band retains 17 better and 13
worse demonstrations. Note that only the rightmost column has a length
AUROC near chance; the two intermediate bands are partial controls and are
reported to make the trend visible, not as controls in themselves.}
\end{table}

Per-timestep normalizations of TED and path length are omitted from these
tables. Unlike SAL, whose raw value grows with the number of spectral bins,
these quantities do not scale linearly with $T$; dividing by $T$ converts
them into inverse-speed proxies that \emph{rise} with skill, which is a
different quantity rather than a duration-free measure of smoothness.

\section{Full downstream training grid}

\begin{table}[h]
\centering
\small
\begin{tabular}{lccc}
\toprule
condition & pool size & best success rate (mean $\pm$ std) & 95\% CI \\
\midrule
oracle (successes only) & 716 & 0.892 $\pm$ 0.048 & [0.866, 0.922] \\
\midrule
no curation             & 1950 & 0.380 $\pm$ 0.063 & [0.344, 0.420] \\
random 50\%             & 1950 & 0.360 $\pm$ 0.061 & [0.326, 0.398] \\
\midrule
SAL                     & 1950 & 0.332 $\pm$ 0.061 & [0.298, 0.370] \\
jerk                    & 1950 & 0.332 $\pm$ 0.041 & [0.306, 0.354] \\
path length             & 1950 & 0.310 $\pm$ 0.030 & [0.292, 0.328] \\
TED                     & 1950 & 0.240 $\pm$ 0.045 & [0.212, 0.266] \\
\bottomrule
\end{tabular}
\caption{Downstream GMM-BC policy success on can/mg, 10 seeds per
condition. Architecture: MLP $1024{\times}1024$, 5 GMM modes, batch 100, 600
epochs of 500 gradient steps, learning rate $10^{-4}$; 50 rollouts
(horizon 400) every 150 epochs, best checkpoint per seed.}
\end{table}

\begin{table}[h]
\centering
\small
\begin{tabular}{lcccc}
\toprule
& \multicolumn{2}{c}{vs.\ random 50\%} & \multicolumn{2}{c}{vs.\ no curation} \\
\cmidrule(lr){2-3}\cmidrule(l){4-5}
metric & $\Delta$ [95\% CI] & $p$ & $\Delta$ [95\% CI] & $p$ \\
\midrule
TED         & $-0.120$ [$-0.166$, $-0.076$] & $\mathbf{0.0001}$ & $-0.140$ [$-0.186$, $-0.096$] & $\mathbf{<0.0001}$ \\
path length & $-0.050$ [$-0.090$, $-0.012$] & 0.037 & $-0.070$ [$-0.112$, $-0.030$] & 0.0071 \\
jerk        & $-0.028$ [$-0.072$, $+0.014$] & 0.25 & $-0.048$ [$-0.092$, $-0.006$] & 0.060 \\
SAL         & $-0.028$ [$-0.080$, $+0.024$] & 0.32 & $-0.048$ [$-0.098$, $+0.002$] & 0.098 \\
\bottomrule
\end{tabular}
\caption{Difference in mean best success rate against each baseline, with
95\% bootstrap CIs on the difference (20{,}000 resamples) and Welch
two-sample $p$-values. Bold marks comparisons surviving Bonferroni
correction across all eight tests ($\alpha = 0.05/8 = 0.006$); only TED
clears it. Random vs.\ no curation: $\Delta = -0.020$, CI
$[-0.072, +0.032]$.}
\end{table}

At 10 seeds, differences below roughly 0.05 in mean success rate are not
resolvable on this grid. A prior 5-seed version of this analysis reported
all four metrics as underperforming both baselines; that pattern held only
for TED and path length once seeds were doubled, and we report the corrected
result.

\begin{table}[h]
\centering
\small
\begin{tabular}{lcc}
\toprule
condition & GMM-BC (10 seeds) & diffusion (5 seeds) \\
\midrule
oracle      & 0.892 $\pm$ 0.048 & 0.544 $\pm$ 0.030 \\
no curation & 0.380 $\pm$ 0.063 & 0.120 $\pm$ 0.014 \\
random 50\% & 0.360 $\pm$ 0.061 & 0.116 $\pm$ 0.017 \\
jerk        & 0.332 $\pm$ 0.041 & 0.212 $\pm$ 0.030 \\
TED         & 0.240 $\pm$ 0.045 & 0.092 $\pm$ 0.030 \\
\bottomrule
\end{tabular}
\caption{Mean best success rate for the five conditions run under both
policy classes.}
\label{tab:crossarch}
\end{table}

\begin{table}[h]
\centering
\small
\begin{tabular}{lcccc}
\toprule
& \multicolumn{2}{c}{vs.\ random 50\%} & \multicolumn{2}{c}{vs.\ no curation} \\
\cmidrule(lr){2-3}\cmidrule(l){4-5}
metric & $\Delta$ [95\% CI] & $p$ & $\Delta$ [95\% CI] & $p$ \\
\midrule
jerk & $+0.096$ [$+0.068$, $+0.124$] & $\mathbf{0.0007}$ & $+0.092$ [$+0.064$, $+0.116$] & $\mathbf{0.0011}$ \\
TED  & $-0.024$ [$-0.048$, $+0.004$] & 0.17 & $-0.028$ [$-0.052$, $+0.000$] & 0.11 \\
\bottomrule
\end{tabular}
\caption{Diffusion-policy grid (Table~\ref{tab:crossarch}), same procedure
as the BC tests above: 95\% bootstrap CIs on the difference (20{,}000
resamples), Welch two-sample $p$-values, Bonferroni across the four tests
reported for this architecture ($\alpha = 0.05/4 = 0.0125$). Random and no
curation are indistinguishable ($\Delta = -0.004$, CI $[-0.020, +0.012]$).
Policy: conditional U-Net ($[256,512,1024]$ down-dims, kernel 5, 8 groups),
EMA power 0.75, observation/action/prediction horizons 2/8/16, DDIM with 10
inference steps, batch 256, 600 epochs of 500 gradient steps, lr $10^{-4}$,
50 rollouts (horizon 400) every 150 epochs. At five seeds this grid resolves
differences of roughly 0.054.}
\end{table}

\section{Secondary setting: can/paired}

On robomimic's matched-pair split every metric is inverted against success
(TED 0.331, SAL 0.398, jerk 0.368, path length 0.029), but the setting
carries an opposing length confound: successful demonstrations are
systematically longer (length-alone AUROC 0.194) and path length correlates
$+0.86$ with duration. Its matched-pair construction also concentrates the
discriminating signal into a handful of terminal timesteps, since each pair
shares a near-identical prefix and diverges only at the final placement.
Neither property is representative of the mixed-quality pools curation
metrics are deployed on, so we report it for completeness and treat can/mg
as the primary success-label evidence.

\section{Qualitative mechanism}

Contrasting a successful and a failed can/mg demonstration: both contain one
high-deviation segment under TED's contact partition, but they differ in
placement and extent. In the success the high-residual segment spans
timesteps 12--96 (84 of 150 steps, mean envelope residual 0.259) and
coincides with the descend-and-grasp. In the failure the only comparable
deviation is the initial free-space approach (0--24, residual 0.213); the
remaining 126 steps are a flat tail in which the end effector holds height
while the gripper cycles without completing a placement, every segment
scoring below 0.014. Because TED averages deviation over the DTW alignment
path, a committed motion spanning half the episode raises the score while a
short approach transient is diluted. TED rates the failure as the cleaner
trajectory (0.042 vs. 0.144). Per-demonstration TED scores inspected:
successes 0.144 and 0.079; failures 0.042, 0.042, 0.017.

\section{Reproducibility}

All data is public robomimic. Metric implementations, scoring scripts, band
analysis, training configurations, and scored CSVs will be released on
publication. TED is reimplemented from the published appendix
($w_{\mathrm{ori}} = 0.25$, contact partition at gripper open/close
transitions, greedy B\'ezier envelope per segment, length-normalized DTW
distance); SAL follows its published formula. Control rate 20\,Hz,
$\Delta t = 0.05$\,s.

\end{document}
